\section{Pre-registered replication on untouched sets}\label{sec:replication}

The main results are descriptive (\S\ref{sec:limits}, Appendix~\ref{app:deviations}). To test the speed claim
once, we ran a replication on two sets that no earlier step used, with a corrected harness and rules written and
committed before any replication data existed (\texttt{REPLICATION\_PLAN.md}); two independent pre-run reviews of
the harness and the plan were applied before the first timing run.

\paragraph{Design.} CQADupStack (all 12 forums merged; 457,199 documents; 13,145 queries, split by forum) and an MS
MARCO subsample of 1,000,000 passages (every passage judged for a dev query plus a seeded random sample; 6,980
queries), embedded with bge-base-en-v1.5. One build per rival family (the build most often frozen for that family in
the main study), every effort-type knob on the same 11-value grid as our stage-1 size, and our cascade only at
configurations the product itself runs. Dev selection at $\fr = 0.99$ used the median of three timed rounds; the rival
of each test is the rival fastest on dev. Rival timers compute only the query forms the family uses, each job
pre-reads only the files it opens, and a job with more than 1,000 major page faults in its timed pass runs once more
(invalid if it fails again). On held-out queries, each CPU type ran 6 fresh instances, each with 2 blocks that hold
ours and every frozen rival point in a seeded random order. The latency test is a one-sided $t$-test over the 6
instance means of log(rival / ours); the quality test is non-inferiority with a relative margin of 0.5\% of exact
quality, bootstrapped jointly over queries; each family of 6 tests is Holm-corrected at 0.025; a set--CPU pair counts
as ``faster at 99\% quality'' only if both tests pass.

\paragraph{Results.} Table~\ref{tab:rep} gives the six pairs. The latency test passes in \nm{rep-lat-pass} of
\nm{rep-n}, the quality test in \nm{rep-q-pass} of \nm{rep-n}, and both in \nm{rep-joint} of \nm{rep-n}. The
geometric mean ratio over instances is \nm{rep-gm} (from \nm{rep-min} to \nm{rep-max}). \nm{rep-reruns} jobs were
rerun by the fault rule and \nm{rep-dropped} blocks were dropped. At the 99th percentile the ratio is \nm{rep-p99};
the native calls are \nm{rep-native}\% of our timed query.

The quality test separates the two sets. On MS MARCO 1M it passes on all three CPUs, so the joint claim (faster and
not worse than the dev-chosen rival by more than 0.5\% of exact quality) holds there on every CPU. On CQADupStack it
fails on all three CPUs: the mean difference to the rival is between $-0.0010$ and $0.0000$ nDCG@10, but its lower
bound falls just below the margin of 0.0021, and both our point and the rival's meet $\fr = 0.99$ only by the point
estimate, not by the bound. The latency advantage on CQADupStack is therefore shown, but quality non-inferiority at
this margin is not; we do not claim it.

\begin{table}[h]\centering\scriptsize
\caption{Replication, held-out, $\fr = 0.99$ (pre-registered tests). Ratio: geometric mean over the 6 instances of
rival / ours (paired bootstrap interval over instances). Latency: one-sided $t$-test, Holm over 6. Quality: mean of
per-query nDCG@10 (ours $-$ rival) and the non-inferiority test at 0.5\% of exact quality, Holm over 6.}\label{tab:rep}
\setlength{\tabcolsep}{3pt}
\begin{tabular}{lllrrrrrlrrl}\toprule
CPU & Set & Rival & Inst. & Ours \us{} & Rival \us{} & Ratio & $p$ (lat.) & Holm & Mean diff & $p$ (qual.) & Holm\\\midrule
\inputtable{tables/rep.tex}
\bottomrule\end{tabular}
\end{table}
